\documentclass[a4paper,12pt]{article}\usepackage[]{graphicx}\usepackage[]{color}
%% maxwidth is the original width if it is less than linewidth
%% otherwise use linewidth (to make sure the graphics do not exceed the margin)
\makeatletter
\def\maxwidth{ %
  \ifdim\Gin@nat@width>\linewidth
    \linewidth
  \else
    \Gin@nat@width
  \fi
}
\makeatother

\definecolor{fgcolor}{rgb}{0.345, 0.345, 0.345}
\newcommand{\hlnum}[1]{\textcolor[rgb]{0.686,0.059,0.569}{#1}}%
\newcommand{\hlstr}[1]{\textcolor[rgb]{0.192,0.494,0.8}{#1}}%
\newcommand{\hlcom}[1]{\textcolor[rgb]{0.678,0.584,0.686}{\textit{#1}}}%
\newcommand{\hlopt}[1]{\textcolor[rgb]{0,0,0}{#1}}%
\newcommand{\hlstd}[1]{\textcolor[rgb]{0.345,0.345,0.345}{#1}}%
\newcommand{\hlkwa}[1]{\textcolor[rgb]{0.161,0.373,0.58}{\textbf{#1}}}%
\newcommand{\hlkwb}[1]{\textcolor[rgb]{0.69,0.353,0.396}{#1}}%
\newcommand{\hlkwc}[1]{\textcolor[rgb]{0.333,0.667,0.333}{#1}}%
\newcommand{\hlkwd}[1]{\textcolor[rgb]{0.737,0.353,0.396}{\textbf{#1}}}%
\let\hlipl\hlkwb

\usepackage{framed}
\makeatletter
\newenvironment{kframe}{%
 \def\at@end@of@kframe{}%
 \ifinner\ifhmode%
  \def\at@end@of@kframe{\end{minipage}}%
  \begin{minipage}{\columnwidth}%
 \fi\fi%
 \def\FrameCommand##1{\hskip\@totalleftmargin \hskip-\fboxsep
 \colorbox{shadecolor}{##1}\hskip-\fboxsep
     % There is no \\@totalrightmargin, so:
     \hskip-\linewidth \hskip-\@totalleftmargin \hskip\columnwidth}%
 \MakeFramed {\advance\hsize-\width
   \@totalleftmargin\z@ \linewidth\hsize
   \@setminipage}}%
 {\par\unskip\endMakeFramed%
 \at@end@of@kframe}
\makeatother

\definecolor{shadecolor}{rgb}{.97, .97, .97}
\definecolor{messagecolor}{rgb}{0, 0, 0}
\definecolor{warningcolor}{rgb}{1, 0, 1}
\definecolor{errorcolor}{rgb}{1, 0, 0}
\newenvironment{knitrout}{}{} % an empty environment to be redefined in TeX

\usepackage{alltt}
\usepackage[applemac]{inputenc}
\usepackage[OT1,T1]{fontenc}
% \usepackage{times}
\usepackage{setspace}
\usepackage{lscape}
\usepackage{amsmath,amssymb}
\usepackage{fancyhdr}
\usepackage{fancyvrb,moreverb,boxedminipage}
\usepackage{multirow}
\usepackage{float}
\usepackage{hyperref}
\usepackage{listings}
\usepackage{verbatim}
\usepackage{enumitem}

\setlength{\textwidth}{160mm}
\setlength{\textheight}{230mm}
\setlength{\oddsidemargin}{-5mm}
\setlength{\topmargin}{-10mm}

\textwidth=6in
\textheight=8in
\topmargin=0.0in
\oddsidemargin=0in 
\baselineskip=18pt
\headheight=2cm
\setcounter{MaxMatrixCols}{10}


\newcommand{\ns}{\!\!\!}
\newcommand{\e}{\varepsilon}
\newcommand{\be}{\beta}
\newcommand{\s}{\sigma}
\newcommand{\vv}{\vspace{.5cm}}
\newcommand{\beps}{\boldsymbol{\epsilon}}
\newcommand{\sumin}{\sum_{i=1}^n}
\newcommand{\bg}[1]{\mbox{\boldmath{$#1$}}}
\DeclareMathOperator{\plim}{plim}
\DefineVerbatimEnvironment{code}{Verbatim}{fontsize=\small}
\DefineVerbatimEnvironment{example}{Verbatim}{fontsize=\small}

\def\by{\mathbf{y}}
\def\bx{\mathbf{x}}
\def\ba{\mathbf{a}}
\def\bb{\mathbf{b}}
\def\be{\mathbf{e}}
\def\bu{\mathbf{u}}
\def\bv{\mathbf{v}}
\def\bz{\mathbf{z}}

%===========
%Greek letter vectors
%===========
\def\bbeta{\mbox{\boldmath $\beta$}}
\def\bgamma{\mbox{\boldmath $\gamma$}}
\def\bvareps{\mbox{\boldmath $\varepsilon$}}
\def\bleta{\mbox{\boldmath $\eta$}}
\def\bmu{\mbox{\boldmath $\mu$}}
\def\btheta{\mbox{\boldmath $\theta$}}
\def\bsigma{\mbox{\boldmath $\sigma$}}
\def\blambgam{\mbox{\boldmath $\lambda_{\gamma}$}}
\def\blambbet{\mbox{\boldmath $\lambda_{\beta}$}}
\def\blambda{\mbox{\boldmath $\lambda$}}

%===========
%Matrices
%===========
\def\bX{\mathbf{X}}
\def\bD{\mathbf{D}(\bsigma)}
\def\bV{\mathbf{V}(\bsigma)}
\def\bVinv{\mathbf{V}^{-1}(\bsigma)}
\def\bR{\mathbf{R}(\bsigma)}
\def\bA{\mathbf{A}}
\def\bB{\mathbf{B}}
\def\bZ{\mathbf{Z}}
\def\bIn{\mathbf{I_{N}}}
\def\bIm{\mathbf{I_{m}}}
\def\bIni{\mathbf{I_{n}}}
\def\SSW{\mathrm{SSW}}
\def\SSB{\mathrm{SSB}}
\def\Normal{\mathcal{N}\!}

%===========
%Parentheses, etc.
%===========
\def\op{\left(}
\def\cp{\right)}
\def\oc{\left[}
\def\cc{\right]}
\def\ok{\left\{}
\def\ck{\right\}}

%===========
%Statistical functions and real numbers.
%===========
\def\exE{\mathbb{E}}
\def\Real{\mathbb{R}}
\def\Var{\mathbb{V}ar}
\def\Proba{\mathbb{P}}
\def\Cov{\mathbb{C}ov}

%===========
%Listings and hyperlink settings
%===========
\hypersetup{colorlinks=true,linkcolor=blue,citecolor=blue, urlcolor=blue}
\urlstyle{tt}

%===========
%Special Environments
%===========
\newenvironment{exercices}{\newcounter{moncompteur}\begin{list}
   {\bfseries Ex. \arabic{moncompteur}}{\usecounter{moncompteur}
     \setlength{\labelwidth}{2.6em}\setlength{\leftmargin}{3.1em}}}{\end{list}}
\renewcommand{\labelenumi}{\bfseries\alph{enumi})}
\newcommand{\splus}[3][0]{\texttt{>~#2}\ \ \ \ \textit{#3}\\[#1mm]}

\newenvironment{exo}
{\begin{center}\setlength{\textwidth}{140mm} \underline{Exercise:} \begin{minipage}[t]{120mm}}
{\end{minipage}\setlength{\textwidth}{160mm}\end{center}}


\usepackage{verbatim}

\lstset{ %
  language=R, 
  basicstyle=\ttfamily,
  backgroundcolor=\color{lightgray},   % choose the background color; you must add \usepackage{color} or \usepackage{xcolor}
  xleftmargin= 10 pt,
  xrightmargin= 10 pt,
  escapeinside={\%*}{*)},          % if you want to add LaTeX within your code
  frame=none,                      % adds a frame around the code
  numbers=none,                    % where to put the line-numbers; possible values are (none, left, right)
  numbersep=5pt,                   % how far the line-numbers are from the code
  numberstyle=\tiny\color{mygray}, % the style that is used for the line-numbers
  showspaces=false,                % show spaces everywhere adding particular underscores; it overrides 'showstringspaces'
}

\pagestyle{fancy}


\definecolor{mygreen}{rgb}{0,0.6,0}
\definecolor{mygray}{rgb}{0.5,0.5,0.5}
\definecolor{mymauve}{rgb}{0.58,0,0.82}
\definecolor{lightgray}{gray}{0.95}
\IfFileExists{upquote.sty}{\usepackage{upquote}}{}
\begin{document}


\lhead{University of Geneva\\ \textbf{GLAM}\\
Pr. Eva Cantoni}
\rhead{GSEM\\ Spring 2022 \\ \textbf{Homework 11}}

\begin{center}
\bigskip

\bigskip

\bigskip

\bigskip

\bigskip

\bigskip

\fbox{{\Large Three full data analyses using \texttt{gam()}}}

\bigskip

\bigskip

\bigskip
\end{center}

%%%%%%%%%%%%%%%%%%%%%%%%%%%%%%%%%%%%%%%%%%%%%%%%%%%%%%%%%%%%%%%%%%%%%%%%%%%%



% \noindent \textsc{University of Geneva} \hfill S411001 SE GLAM\\
% Prof. Eva Cantoni \hfill Spring semester 2015\\
% \noindent
% \makebox[\linewidth]{\rule{\textwidth}{0.4pt}}\\[1.5ex]
% \textbf{} \hfill \textbf{Homework 11: Three full data analyses using \texttt{gam()}}\\
% \makebox[\linewidth]{\rule{\textwidth}{0.4pt}}\\
% 
% \noindent Course website: \url{https://chamilo.unige.ch/home/courses/S411001/}

\setlength{\parskip}{1ex plus 0.5ex minus 0.2ex}
\parindent=0cm
\setlength{\parskip}{1ex plus 0.5ex minus 0.2ex}
\linespread{1.1}
%\sloppy

\section{\underline{Exercise 1}}

To illustrate the use of a GAM on a logistic example, we consider the vasoconstriction data set. The data consist of 39 binary responses denoting the presence or absence (coded 1/0 respectively) of vasoconstriction (\verb"vasoconstr") on the fingers' skin after inspiration of a volume of air (\verb"vol") at a rate reported in the variable called \verb"rate".
\begin{enumerate}[label=\arabic*.]
\item{Consider first taking the logarithm of both predictors, as suggested by many authors that have analyzed this data set.
\begin{enumerate}
\item{Estimate a logistic additive model. Comment the plot of the estimated functions and the residuals.}
\item{Based on the AIC criteria, would a GLM be sufficient~?}
\item{Is the \verb"rate" variable useful in our model?}\end{enumerate}}
\item{Fit now the same model but without transforming the covariates.}
\item{Estimate finally a bivariate surface and compare it with the two previous models.}
\end{enumerate}

\section{\underline{Exercise 2}}

Data were collected on 83 patients who underwent corrective spinal surgery (Bell et al., 1989). The objective of the study is to determine the important risk factors for the curvature of the spine (at least 40 degrees with respect to the vertical), following surgery. The \verb"Kyphosis" variable describes the presence/absence of the curvature coded respectively 1/0. The risk factors are \verb"Age" in months, the starting (\verb"Start") vertebrae level of the surgery (range 1-12 corresponds to the thoracic vertebrae, while 13-17 are the lumbar ones) and the \verb"Number" of vertebrae levels involved.
\begin{enumerate}
\item{Source the \verb"PairsLab" function. Produce the pairwise scatter plots of all three predictors, with the responses contained in the \verb"Kyphosis" variable replacing the plot symbols (you specify this through the \texttt{label} argument of the \texttt{PairsLab} function). What can you observe concerning the regions of 1s and 0s~?}
\item{Fit an additive logistic model using all three predictors.}
\item{%If you look carefully at the patterns in the partial residuals plots for \verb"Age" %and \verb"Start", what do you remark
   %(hint: known as pure-region effects)~?
   Are there any values in the \verb"Age" and \verb"Number" variables which seem to have a large effect on the function estimates (with respect to the pointwise standard-error bands)?}
\item{Do the same fit with those observations removed. What happens~?}
\item{Are all the three predictors important~? Enquire if you should conserve $f(Number)$ as explanatory variable. Then, estimate your final model.}
\end{enumerate}


\section{\underline{Exercise 3}}

We would like to model the number of cigarettes smoked per day among smokers according to some available covariates\footnote{Data from Wooldridge, J. M. (2002). Econometric Analysis of Cross Section and Panel Data. MA: MIT Press}. We have data on 310 individuals all living in the USA, but in different states. Due to the different state legislatures, cigarettes prices and smoking restrictions vary from one individual to the other. Our data feature the following variables:
\begin{itemize}
\item{the response \texttt{cigs}, the number of cigarettes smoked per day;}
\item{\texttt{cigpric}, the price (in ten cents) for a cigarette pack in the individual's state;}
\item{\texttt{income} in US dollars;}
\item{the dummy \texttt{restaurn} indicating whether there are restaurant smoking restrictions in the state or not (\texttt{1} = yes, \texttt{0} = no);}
\item{the dummy \texttt{white} coding for the individual's skin color (\texttt{1} = white, \texttt{0} = non-white);}
\item{\texttt{educ}, the number of years of education;}
\item{and \texttt{age} given in years.}
\end{itemize}

\begin{enumerate}
\item{Comment the distribution of the response.}
\item{Fit a GAM to these data. Which variables are you going to keep in the parametric part~?}
\item{Interpret the overall effect of age on the number of cigarettes smoked per day.}
\item{Can we include income in a linear fashion in our model? If not, how should we transform this variable if we still want it in the parametric part~?}
\item{Does the price of a cigarette pack really have an impact on the number cigarettes smoked per day~?}
\end{enumerate}




\end{document}
